\chapter{Đào sâu: Dữ liệu xuyên service --- mỗi service một cơ sở dữ liệu, dữ liệu trùng lặp và outbox}
\label{ch:dd-consistency}

Chương~7 đặt ra quy tắc: mỗi service sở hữu cơ sở dữ liệu của mình, và
con đường duy nhất tới dữ liệu của service khác là API hoặc sự kiện của
nó. Quy tắc dễ phát biểu và tốn kém để sống cùng. Chương đào sâu này đi
theo mẩu dữ liệu duy nhất mà dự án sao chép giữa các service --- tên và
e-mail của một học viên --- từ lúc nó được ghi vào Postgres của
auth-service tới lúc nó xuất hiện trong Postgres của course-service, và
hỏi chuyện gì có thể hỏng trên đường đi. Câu trả lời là bài toán ghi kép
(dual write), và cách sửa là transactional outbox.

Ba câu hỏi sắp xếp chương này:

\begin{enumerate}
  \item \textbf{Mỗi service một cơ sở dữ liệu mua được gì}, và cấm điều
  gì --- join, khóa ngoại, transaction --- mà một monolith coi là hiển
  nhiên?
  \item \textbf{Dữ liệu trùng lặp giữ đúng bằng cách nào}, và cửa sổ
  chính xác mà dự án đang sai hôm nay là gì?
  \item \textbf{``Lưu dòng và phát sự kiện'' có thể nguyên tử được
  không} --- và nếu không, thứ gần nhất với nó là gì?
\end{enumerate}

\section{Vì sao chọn nó, và nó giải quyết gì}

Mở \code{docker-compose.yml}: hai container Postgres~17,
\code{postgres-auth} giữ \code{ts\_auth} và \code{postgres-course} giữ
\code{ts\_course}, trên volume và cổng riêng. Không gì trong
course-service biết địa chỉ của cái đầu. Sự cô lập mua bốn thứ:

\begin{itemize}
  \item \textbf{Quyền sở hữu schema.} Auth-service có thể thêm cột
  \code{totp\_secret}, đổi tên \code{provider\_id}, hay chạy một
  migration Flyway phá hủy mà không cần phối hợp với ai. Các migration
  \code{V3} và \code{V4} ở Chương~7 đã diễn ra đúng như thế.
  \item \textbf{Triển khai độc lập.} Một migration khóa bảng
  \code{users} một phút làm nghẽn đăng nhập, không phải danh sách khóa
  học.
  \item \textbf{Mở rộng và tinh chỉnh độc lập.} Cơ sở dữ liệu khóa học
  sẽ lớn theo bài nộp; cơ sở dữ liệu auth thì không. Chúng có thể nằm
  trên đĩa khác nhau, có connection pool khác nhau, lịch sao lưu khác
  nhau.
  \item \textbf{Chặn cứng sự ghép nối vô tình.} Trong một cơ sở dữ
  liệu dùng chung, câu \code{JOIN users ON student.user\_id = users.id}
  đầu tiên viết dưới áp lực deadline hàn schema của hai đội lại với nhau
  mãi mãi. Ở đây join là bất khả, nên nó không bao giờ được viết.
\end{itemize}

Nó cấm đúng bằng ấy:

\begin{itemize}
  \item \textbf{Không join xuyên service.} ``Liệt kê học viên khóa~7
  kèm e-mail'' không thể là một truy vấn nếu e-mail nằm ở nơi khác.
  \item \textbf{Không khóa ngoại xuyên service.} \code{student.user\_id}
  trong \code{V1\_\_create\_course\_tables.sql} là \code{BIGINT NOT NULL
  UNIQUE} và không gì hơn --- Postgres không thể kiểm tra user~42 tồn
  tại trong một cơ sở dữ liệu nó không nhìn thấy.
  \item \textbf{Không transaction xuyên service.} Một thao tác nghiệp
  vụ phải đổi cả hai cơ sở dữ liệu không thể rollback như một khối.
\end{itemize}

\begin{decision}[sao chép học viên, không gọi lấy user]
Course-service cần tên và e-mail học viên trên mọi danh sách lớp và
trong mọi thông báo bài tập. Có hai thiết kế: gọi auth-service mỗi lần
(thậm chí có sẵn một Feign client cho việc đó,
\code{AuthServiceClient.getUserById}, với fallback trả về \code{null}),
hoặc giữ một bản sao cục bộ trong bảng \code{student}. Dự án chọn bản
sao, và Feign client là code chết --- không class nào trong
course-service gọi nó. Lý do: danh sách lớp được đọc nhiều hơn rất
nhiều so với số lần tên thay đổi; một lời gọi đồng bộ cho mỗi dòng sẽ
khiến tính sẵn sàng của course-service là tích của tính sẵn sàng
auth-service (Chương~6); và bản sao cục bộ có thể index, join, phân
trang bằng SQL thuần. Cái giá là bản sao có thể \emph{cũ}, và phần còn
lại của chương nói về cũ đến đâu, và làm sao chặn trên nó. Quyết định
lật lại khi dữ liệu thay đổi thường xuyên hoặc phải là nguồn tin cậy
lúc đọc --- số dư tài khoản, một quyền hạn --- nơi một lời gọi, hoặc
cache với vô hiệu hóa tường minh, mới là thiết kế trung thực.
\end{decision}

\section{Một học viên tới course-service bằng cách nào}

Lần theo đường đi thật. Nó bắt đầu ở \code{AuthService} của
auth-service, được đánh dấu \code{@Transactional} ở mức class:

\begin{lstlisting}[language=Java]
public void activate(String code) {
    User user = userRepository.findByActivationCode(code)
            .orElseThrow(() ->
                new IllegalArgumentException("Invalid activation code"));

    user.setActivated(true);
    user.setActivationCode(null);
    userRepository.save(user);                       (*@\co{1}@*)

    userActivatedPublisher.publish(user);            (*@\co{2}@*)
}
\end{lstlisting}

\code{UserActivatedPublisher} dựng một record \code{UserActivatedEvent}
--- id user, e-mail, họ tên, vai trò --- và gọi
\code{kafkaTemplate.send("ts.user.activated", event.userId(), event)}.
Ở phía kia, \code{UserActivatedListener} của course-service:

\begin{lstlisting}[language=Java]
@KafkaListener(topics = "ts.user.activated", groupId = "course-group")
public void onUserActivated(UserActivatedEvent event) {
    if (!ROLE_STUDENT.equalsIgnoreCase(event.role())) {
        return;
    }
    Long userId = Long.parseLong(event.userId());
    if (studentRepository.findByUserId(userId).isPresent()) { (*@\co{3}@*)
        return;
    }
    Student student = Student.builder()
            .userId(userId)
            .firstName(event.firstName())
            .lastName(event.lastName())
            .email(event.email())
            .build();
    studentRepository.save(student);                       (*@\co{4}@*)
}
\end{lstlisting}

\begin{description}
  \item[\co{1}] Dòng đang bẩn bên trong một transaction chưa commit.
  \code{save()} trên một entity đã được quản lý là no-op cho tới khi
  flush.
  \item[\co{2}] Sự kiện được trao cho Kafka producer \emph{bên trong}
  cùng transaction, trước commit ở lúc thoát method. Ghi nhớ dòng này;
  nó là toàn bộ vấn đề.
  \item[\co{3}] Listener lũy đẳng (idempotent) theo khóa tự nhiên: giao
  lại lần hai cùng sự kiện thấy dòng đã có và trả về. Đây là lý do dự án
  sống sót qua giao nhận ít-nhất-một-lần (Chương~34) mà không cần bảng
  khử trùng --- ở đây.
  \item[\co{4}] Bản sao. Từ giờ, course-service trả lời mọi truy vấn
  danh sách lớp từ dòng này và không bao giờ hỏi auth-service nữa.
\end{description}

Cửa sổ nhất quán vì thế là: từ commit ở auth-service tới commit ở
course-service. Trên một stack yên tĩnh đó là mili giây. Với
notification-service và course-service cùng nằm trên đường consumer,
và broker trên cùng cái máy hai nhân (Chương~25), nó có thể là vài
giây. Một giáo viên ghi danh học viên ``ngay lập tức'' sau khi học viên
kích hoạt có thể nhận \code{404 Student not found} từ
\code{EnrollmentService.enroll} --- hành vi đúng, trải nghiệm bất ngờ.

\subsection{Thứ hoàn toàn không được lan truyền}

Giờ tìm sự kiện thứ hai. Không có. \code{ts.user.activated} bắn một
lần, lúc kích hoạt. Một user sau đó đổi tên hay e-mail trong
auth-service không bao giờ được phản ánh vào \code{student}; ghi chú
``future: Kafka event propagates e-mail changes'' trong
\code{DATABASE-DESIGN.md} là lời thừa nhận. Thông báo bài tập sẽ tiếp
tục đi tới địa chỉ cũ. Đây là trạng thái trung thực của nhất quán cuối
cùng trong dự án: \emph{cuối cùng} cho việc tạo, \emph{không bao giờ}
cho việc cập nhật. Mục~\ref{sec:enhancing-outbox} thêm sự kiện còn
thiếu.

\section{Ưu và nhược, thẳng thắn}

\begin{center}
\begin{tabular}{@{}p{0.3\linewidth}p{0.3\linewidth}p{0.3\linewidth}@{}}
\toprule
 & Cơ sở dữ liệu dùng chung & Mỗi service một CSDL (ở đây) \\
\midrule
Đọc xuyên entity & một \code{JOIN} & bản sao cục bộ, hoặc hai lời gọi \\
Toàn vẹn & khóa ngoại & code ứng dụng + sự kiện \\
Ghi nhiều entity & một transaction & saga hoặc outbox \\
Đổi schema & phối hợp mọi chủ sở hữu & chủ sở hữu tự quyết \\
Bán kính sự cố & một CSDL, mọi service & một service \\
Độ cũ & không & vài giây (tạo), không giới hạn (cập nhật) hôm nay \\
\bottomrule
\end{tabular}
\end{center}

Điểm yếu thật của dự án, gọi tên:

\begin{itemize}
  \item Sự kiện cập nhật không tồn tại; bản sao trôi vĩnh viễn.
  \item Lưu-và-phát không nguyên tử (mục kế tiếp).
  \item Fallback của Feign trả \code{null}, nên nếu ai đó nối
  \code{AuthServiceClient} vào, trạng thái circuit mở sẽ hiện ra dưới
  dạng \code{NullPointerException} thay vì một phản hồi suy giảm rõ
  ràng.
  \item Không có đường báo cáo: ``bao nhiêu học viên đã kích hoạt mỗi
  khóa'' cần dữ liệu từ cả hai cơ sở dữ liệu và không có read model nào
  cho nó.
\end{itemize}

\section{Bài toán ghi kép}
\label{sec:dual-write}

Dòng \co{2} ở trên ghi vào hai hệ thống --- Postgres và Kafka --- không
chia sẻ transaction manager. Dù bạn đặt thứ tự thế nào, một trong hai
sự cố là khả dĩ:

\begin{itemize}
  \item \textbf{Commit, rồi mất sự kiện.} \code{send()} là bất đồng
  bộ: nó nối vào một bộ đệm phía client và trả về một future. Method
  thoát, transaction commit, và JVM bị giết (\code{OOMKilled}, một
  rolling update, \code{kubectl delete pod}) trước khi bộ đệm được
  flush. User đã kích hoạt; course-service không bao giờ nghe về nó;
  học viên không bao giờ ghi danh được, và không gì log lỗi.
  \item \textbf{Phát, rồi rollback.} Bộ đệm flush trong vài mili giây;
  broker ack; rồi commit thất bại --- một ràng buộc unique, kết nối
  rớt, deadlock. Course-service tạo học viên cho một user không tồn
  tại. Một dòng \emph{ma}.
\end{itemize}

Cả hai hiếm theo từng request và chắc chắn theo từng năm. Phiên bản
phỏng vấn: ``bạn lưu một đơn hàng và phát \code{OrderCreated}; làm sao
bảo đảm cả hai xảy ra?'' Câu trả lời junior là ``bọc trong try/catch''.
Câu trả lời senior bắt đầu bằng ``không thể, với hai tài nguyên'', rồi
gọi tên outbox.

\begin{pitfall}[đăng ký treo sáu mươi giây, rồi thất bại]
Triệu chứng: khi Kafka chết, \code{POST /api/auth/register} không thất
bại nhanh --- nó treo một phút rồi trả 500. Nguyên nhân: producer phải
lấy metadata topic trước lần gửi đầu; không broker nào với được, nó
chặn trong \code{max.block.ms} (mặc định 60\,000\,ms) rồi ném ngoại lệ.
Ngoại lệ lan ra khỏi \code{register()}, transaction rollback, và user
\emph{không} được đăng ký --- vì broker gửi mail chết. Một phụ thuộc lẽ
ra bất đồng bộ đã thành điểm lỗi đồng bộ với timeout một phút. Outbox
bên dưới biến chuyện này thành ``đã đăng ký; e-mail đến trễ''.
\end{pitfall}

\begin{pitfall}[bản sao sống lâu hơn bản gốc]
Triệu chứng: e-mail bài tập bị trả lại từ một địa chỉ user đã đổi nhiều
tháng trước. Nguyên nhân: \code{student.email} được
\code{UserActivatedListener} ghi một lần và không có sự kiện
\code{ts.user.updated}. Cách sửa: phát ra một sự kiện như thế, và để
listener upsert. Mục~\ref{sec:enhancing-outbox} cho thấy bên phát; thay
đổi ở listener là một
\code{findByUserId().map(update).orElseGet(create)}.
\end{pitfall}

\section{Nâng cấp giải pháp: transactional outbox}
\label{sec:enhancing-outbox}

Ý tưởng là ghi sự kiện vào \emph{cùng cơ sở dữ liệu}, trong \emph{cùng
transaction}, với dòng nghiệp vụ --- và để một tiến trình riêng chép nó
sang Kafka sau đó. Postgres bảo đảm dòng và sự kiện commit cùng nhau
hoặc không gì cả; relay bảo đảm sự kiện cuối cùng tới Kafka; tính lũy
đẳng của consumer bảo đảm relay thử lại không gây hại. Đây là cài đặt
đầy đủ cho auth-service.

Bảng, dưới dạng migration Flyway \code{V5\_\_create\_outbox.sql}:

\begin{lstlisting}[language=SQL]
CREATE TABLE outbox (
    id             UUID PRIMARY KEY,
    aggregate_type VARCHAR(64)  NOT NULL,   -- 'user'
    aggregate_id   VARCHAR(64)  NOT NULL,   -- also the Kafka key
    event_type     VARCHAR(64)  NOT NULL,   -- 'UserActivated'
    topic          VARCHAR(128) NOT NULL,
    payload        JSONB        NOT NULL,
    created_at     TIMESTAMPTZ  NOT NULL DEFAULT now(),
    published_at   TIMESTAMPTZ,
    attempts       INT          NOT NULL DEFAULT 0
);
-- the relay's only query: oldest unpublished first
CREATE INDEX idx_outbox_unpublished
    ON outbox (created_at) WHERE published_at IS NULL;
\end{lstlisting}

Entity cố tình ngu ngốc --- một dòng, không phải domain object:

\begin{lstlisting}[language=Java]
@Entity @Table(name = "outbox")
@Getter @Setter @NoArgsConstructor
public class OutboxEvent {
    @Id private UUID id;
    private String aggregateType;
    private String aggregateId;
    private String eventType;
    private String topic;
    @JdbcTypeCode(SqlTypes.JSON)
    private String payload;                       // serialized event
    private Instant createdAt;
    private Instant publishedAt;
    private int attempts;
}

public interface OutboxRepository extends JpaRepository<OutboxEvent, UUID> {
    // SKIP LOCKED: two relay instances never pick the same row
    @Query(value = """
        SELECT * FROM outbox WHERE published_at IS NULL
        ORDER BY created_at LIMIT :n FOR UPDATE SKIP LOCKED""",
        nativeQuery = true)
    List<OutboxEvent> lockUnpublished(@Param("n") int n);
}
\end{lstlisting}

Publisher không còn đụng tới Kafka. Nó ghi một dòng, và vì được gọi từ
bên trong transaction của \code{AuthService}, dòng đó commit cùng với
user:

\begin{lstlisting}[language=Java]
@Component
@RequiredArgsConstructor
public class UserActivatedPublisher {

    private final OutboxRepository outbox;
    private final ObjectMapper mapper;

    public void publish(User user) {
        UserActivatedEvent event = new UserActivatedEvent(
                String.valueOf(user.getId()), user.getEmail(),
                user.getFirstName(), user.getLastName(),
                user.getRole().name());
        OutboxEvent row = new OutboxEvent();
        row.setId(UUID.randomUUID());               (*@\co{1}@*)
        row.setAggregateType("user");
        row.setAggregateId(event.userId());          (*@\co{2}@*)
        row.setEventType("UserActivated");
        row.setTopic("ts.user.activated");
        row.setPayload(toJson(event));
        row.setCreatedAt(Instant.now());
        outbox.save(row);                            (*@\co{3}@*)
    }

    private String toJson(Object o) {
        try { return mapper.writeValueAsString(o); }
        catch (JsonProcessingException e) {
            throw new IllegalStateException(e);
        }
    }
}
\end{lstlisting}

Relay chạy mỗi giây, lấy một lô đã khóa, gửi từng dòng \emph{đồng bộ},
và đánh dấu:

\begin{lstlisting}[language=Java]
@Component
@RequiredArgsConstructor
@Slf4j
public class OutboxRelay {

    private final OutboxRepository outbox;
    private final KafkaTemplate<String, String> kafka;

    @Scheduled(fixedDelay = 1000)
    @Transactional                                   (*@\co{4}@*)
    public void relay() {
        for (OutboxEvent row : outbox.lockUnpublished(100)) {
            ProducerRecord<String, String> rec = new ProducerRecord<>(
                    row.getTopic(), row.getAggregateId(), row.getPayload());
            rec.headers().add("eventId",
                    row.getId().toString().getBytes(UTF_8));    (*@\co{5}@*)
            rec.headers().add("eventType",
                    row.getEventType().getBytes(UTF_8));
            try {
                kafka.send(rec).get(10, TimeUnit.SECONDS);      (*@\co{6}@*)
                row.setPublishedAt(Instant.now());
            } catch (Exception e) {
                row.setAttempts(row.getAttempts() + 1);
                log.warn("outbox {} attempt {} failed: {}",
                        row.getId(), row.getAttempts(), e.toString());
                break;                                          (*@\co{7}@*)
            }
        }
    }
}
\end{lstlisting}

\begin{description}
  \item[\co{1}] Id sự kiện được đúc \emph{ở đây}, một lần, và sống qua
  mọi lần thử lại. Nó là khóa khử trùng của consumer (Chương~34).
  \item[\co{2}] Id aggregate trở thành key Kafka, nên mọi sự kiện về
  một user rơi vào một partition, theo thứ tự.
  \item[\co{3}] Một \code{INSERT} thuần, trong transaction của bên gọi.
  Nếu \code{activate()} rollback, dòng này cũng vậy: không sự kiện ma.
  Nếu commit, dòng đã bền vững trước khi Kafka dính líu: không mất sự
  kiện.
  \item[\co{4}] Transaction của relay giữ các khóa \code{FOR UPDATE SKIP
  LOCKED}, nên replica auth-service thứ hai chạy cùng scheduler lấy
  những dòng khác.
  \item[\co{5}] Header mang danh tính và loại mà không làm bẩn payload
  JSON mà các class DTO giải tuần tự.
  \item[\co{6}] Nơi \emph{duy nhất} trong auth-service chờ Kafka, và nó
  là một luồng nền. Request HTTP không bao giờ chặn trên broker nữa.
  \item[\co{7}] Khi thất bại thì dừng lô, giữ thứ tự (dòng~2 không được
  phát trước dòng~1 cho cùng key), và để nhịp kế tiếp thử lại.
  \code{attempts} là thứ bạn đặt cảnh báo.
\end{description}

Thêm \code{@EnableScheduling} vào \code{AuthServiceApplication} và đổi
\code{KafkaTemplate} thành \code{<String, String>}, vì payload đã là
JSON. Consumer thấy đúng thứ nó từng thấy.

Hai dòng sổ sách hoàn tất nó: một lệnh hằng đêm
\code{DELETE FROM outbox WHERE published\_at < now() - interval '7 days'}
giữ bảng nhỏ, và cùng hình dạng \code{UserActivatedPublisher} cho bạn
\code{ts.user.updated} miễn phí --- gọi nó từ bất cứ đâu lưu thay đổi
hồ sơ, với \code{eventType = "UserUpdated"}, và để
\code{UserActivatedListener} upsert thay vì insert.

\begin{handson}[xem outbox hấp thụ một sự cố broker]
Với outbox đã có, dừng Kafka, kích hoạt một user, và nhìn:
\begin{lstlisting}[language=bash]
$ docker compose stop kafka
$ curl -s -X POST localhost:8080/api/auth/activate?code=$CODE
{"message":"Account activated"}          # instant, no 60s hang
$ docker compose exec postgres-auth psql -U postgres ts_auth -c \
  "select event_type, attempts, published_at from outbox"
 event_type    | attempts | published_at
 UserActivated |        3 |
$ docker compose start kafka
# ...a few seconds later
 UserActivated |        3 | 2026-09-07 03:04:11.20+00
\end{lstlisting}
Rồi xem log của course-service tìm \code{Provisioned student row}. Việc
kích hoạt chưa bao giờ gặp rủi ro, và học viên xuất hiện ngay khi
broker xuất hiện.
\end{handson}

\subsection{Change data capture: relay cấp công nghiệp}

Relay kiểu polling đơn giản và thêm tối đa một giây độ trễ cộng một
truy vấn mỗi giây mỗi replica. Ở quy mô lớn relay được thay bằng
\emph{change data capture}: Debezium đọc đuôi write-ahead log của
Postgres, thấy \code{INSERT INTO outbox}, và tự phát nó lên Kafka,
không cần scheduler ứng dụng và độ trễ dưới 100\,ms. Code ứng dụng y
hệt --- ghi dòng --- đó là lý do nên áp dụng outbox sớm: relay là phần
duy nhất thay đổi.

\section{Saga: khi hai service đều phải thay đổi}

Ghi danh hôm nay là cục bộ --- \code{course}, \code{student} và
\code{enrollment} đều nằm trong \code{ts\_course}, nên
\code{EnrollmentService.enroll} là một transaction bình thường. Lần ghi
xuyên service thực sự đầu tiên mà dự án sẽ gặp là \emph{xóa tài
khoản}: auth-service phải xóa user; course-service phải xóa học viên,
các ghi danh, bài nộp, và file bài nộp trong MinIO (Chương~9). Không
có transaction nào trải qua từng ấy.

Một \emph{saga} là chuỗi các transaction cục bộ trong đó mỗi bước phát
một sự kiện kích hoạt bước kế, và mỗi bước có một hành động \emph{bù}
(compensating) nếu bước sau thất bại:

\begin{enumerate}
  \item auth-service đánh dấu user \code{DELETING} (chưa xóa) và ghi
  \code{UserDeletionRequested} vào outbox.
  \item course-service xóa dữ liệu và file của học viên, rồi phát
  \code{StudentDataDeleted} --- hoặc, khi thất bại,
  \code{StudentDataDeletionFailed}.
  \item auth-service, khi nhận sự kiện thành công, xóa hẳn user; khi
  nhận sự kiện thất bại, bù bằng cách khôi phục \code{ACTIVE} và báo
  cho vận hành.
\end{enumerate}

Đây là \emph{choreography}: không có điều phối viên, mỗi service phản
ứng với sự kiện. Nó hợp với hai hay ba bước. Quá số đó --- xóa xuyên
năm service, với ràng buộc thứ tự --- đồ thị sự kiện trở nên không đọc
nổi, và các đội chuyển sang \emph{orchestration}: một thành phần (saga
orchestrator, hoặc một workflow engine như Temporal) giữ máy trạng thái
và gửi \emph{lệnh} tới từng service. Đánh đổi là thêm một service để
vận hành lấy một luồng bạn thực sự đọc được.

Hai tính chất mọi saga phải có, và người phỏng vấn kiểm tra:
\emph{khóa ngữ nghĩa} --- trạng thái \code{DELETING} ngăn user đăng nhập
trong vài giây saga chạy --- và \emph{bước lũy đẳng}, để giao nhận
ít-nhất-một-lần của Chương~34 không thể xóa hai lần hay bù hai lần.

\section{Đọc xuyên service}

Khoảng trống còn lại là báo cáo. ``Học viên đã kích hoạt mỗi khóa, kèm
lần đăng nhập cuối'' cần dữ liệu auth và dữ liệu khóa học, và không
service nào có cả hai. Ba câu trả lời, theo chi phí tăng dần:

\begin{itemize}
  \item \textbf{Ghép API.} Gateway hoặc một backend-for-frontend nhỏ gọi
  cả hai service và join trong bộ nhớ. Ổn cho một trang, sai cho một
  báo cáo hằng đêm trên mọi khóa học.
  \item \textbf{Một read model.} Một cơ sở dữ liệu \code{reporting}
  riêng đăng ký \emph{mọi} sự kiện miền và duy trì các bảng phi chuẩn
  hóa theo hình dạng truy vấn. Đây là CQRS: ghi đi tới các service sở
  hữu, đọc đi tới một mô hình dựng từ sự kiện của chúng. Outbox ở trên
  là điều kiện tiên quyết --- không có sự kiện đáng tin, read model trôi
  âm thầm.
  \item \textbf{Một kho dữ liệu.} CDC từ mọi cơ sở dữ liệu service vào
  một kho phân tích. Read model cho những người không phải bạn.
\end{itemize}

\section{Chuyện gì gãy lúc 3 giờ sáng}

\begin{itemize}
  \item \textbf{postgres-course chết; user vẫn kích hoạt.} Không có
  outbox, không gì mất ở phía auth, nhưng \code{UserActivatedListener}
  ném ngoại lệ ở \code{save()}. Error handler mặc định của Spring thử
  lại mười lần liên tiếp, rồi \emph{log và bỏ qua} record (Chương~34).
  Sự kiện đã tiêu thụ, học viên không bao giờ được tạo, và dòng log là
  dấu vết duy nhất. Có outbox thì phía auth an toàn, nhưng phía consumer
  vẫn cần topic thử lại --- outbox sửa producer, không sửa consumer.
  \item \textbf{Một sự kiện bị phát lại.} Ai đó reset offset của
  \code{course-group} để phục hồi sau sự cố trên. Mọi
  \code{ts.user.activated} kể từ đầu retention được giao lại.
  \code{findByUserId().isPresent()} biến từng cái thành no-op. Lũy đẳng
  theo khóa tự nhiên là thứ biến phát lại thành thao tác thường ngày
  thay vì sự cố.
  \item \textbf{Thứ tự.} \code{UserActivated} và một
  \code{UserUpdated} tương lai cho cùng user được key theo
  \code{userId} nên chung partition; cập nhật không thể vượt mặt kích
  hoạt. Sự kiện cho các user \emph{khác nhau} có thể xen kẽ tự do, ổn
  thôi --- không có bất biến xuyên user. Đổi key, hoặc gửi từ hai
  producer khác nhau, và bảo đảm này biến mất.
  \item \textbf{Relay kẹt.} Một dòng độc --- payload broker từ chối vì
  kích thước --- khiến \co{7} break ở mọi nhịp, và các dòng phía sau xếp
  hàng mãi mãi. Cảnh báo trên \code{max(attempts)} và trên
  \code{count(*) where published\_at is null}; sau N lần thử, chuyển
  dòng sang bảng \code{outbox\_dead} và để hàng đợi chảy tiếp.
\end{itemize}

\begin{decision}[khi một cơ sở dữ liệu là lựa chọn tốt hơn]
\code{docs/PLAN-A} mô tả cùng sản phẩm dưới dạng modular monolith: một
ứng dụng Spring Boot, một cơ sở dữ liệu \code{ts}, và
\code{student.user\_id} là khóa ngoại thật. Mọi vấn đề trong chương này
--- ghi kép, trôi dữ liệu, saga, read model --- biến mất ở đó, thay bằng
một \code{JOIN} và một transaction. Đó không phải khoản tiết kiệm nhỏ.
Bố cục microservice (\code{PLAN-B}) chỉ xứng đáng với chi phí khi những
điều bị cấm chính là mục đích: các đội phải triển khai không cần nhau,
dữ liệu phải mở rộng riêng, hoặc tổ chức đủ lớn để một schema dùng
chung \emph{sẽ} bị join chéo trong cơn nóng giận. Với một đội ba người
dự án này cố ý xây quá tay, và chặng đường phía trước ở Chương~27
thẳng thắn về điều đó. Câu trả lời senior trong phỏng vấn không phải
``microservices'' hay ``monolith''; nó là ``modular monolith cho tới
khi một ranh giới có lý do để trở thành ranh giới mạng --- và khi đó,
outbox ngay từ ngày đầu''.
\end{decision}

\section{Ngoài dự án}

Outbox là mẫu đứng sau ``outbox event router'' của Debezium, transactional
messaging của Axon và Eventuate, và các dạng
\code{TransactionalEventPublisher} trong mọi codebase hướng sự kiện
nghiêm túc. Saga kiểu orchestration được sản phẩm hóa bởi Temporal,
Camunda và AWS Step Functions; kiểu choreography đơn giản là ``sự kiện
kèm hành động bù'' và một sơ đồ ai đó phải giữ cho cập nhật. Netflix,
Uber, và mọi ngân hàng bạn kể tên được đều chạy một tổ hợp nào đó của
ba thứ này --- outbox cho độ tin cậy, saga cho ghi nhiều bước, read
model CQRS cho truy vấn vượt ranh giới sở hữu. Điều duy nhất không ai
trong số họ làm là điều dự án đang làm hôm nay: gọi \code{send()} bên
trong một transaction và hy vọng.

\section{Câu hỏi từ phỏng vấn thật}

\subsection*{Q: Bạn cần một báo cáo join user và khóa học, mà chúng nằm ở hai cơ sở dữ liệu. Làm thế nào?}

Trước hết tôi hỏi cần tươi đến đâu và lớn đến đâu. Cho một màn hình,
ghép API: một backend-for-frontend gọi \code{/users?ids=} và
\code{/courses}, join trong bộ nhớ, ổn tới vài trăm dòng. Cho một báo
cáo thật tôi dựng read model: một cơ sở dữ liệu \code{reporting} đăng
ký \code{ts.user.activated}, \code{ts.user.updated} và các sự kiện khóa
học, duy trì một bảng \code{course\_roster} phi chuẩn hóa theo hình
dạng truy vấn. Đó là CQRS --- các service sở hữu giữ phần ghi, read
model trả lời phần đọc xuyên service. Điều kiện tiên quyết là sự kiện
đáng tin, nên tôi đặt outbox vào trước; một read model nuôi bằng những
lời gọi \code{send()} kiểu cố-gắng-hết-sức sẽ trôi và không ai nhận ra
cho tới khi bộ phận tài chính nhận ra. Cho phân tích ở quy mô công ty,
cùng ý tưởng là CDC vào kho dữ liệu. Điều tôi từ chối làm là cấp cho
job báo cáo một read replica của cả hai cơ sở dữ liệu và một join xuyên
cơ sở dữ liệu --- đó là sự ghép nối qua cơ sở dữ liệu dùng chung mà kiến
trúc này tồn tại để ngăn, quay lại bằng cửa hông.

\subsection*{Q: Hai service phải cùng thành công hoặc không service nào. Làm thế nào, không dùng two-phase commit?}

Một saga. Mỗi service làm một transaction cục bộ và phát sự kiện; bước
kế phản ứng; mỗi bước có hành động bù. Trong dự án, xóa tài khoản:
auth-service đặt \code{DELETING} và phát \code{UserDeletionRequested};
course-service xóa dữ liệu và phát \code{StudentDataDeleted} hoặc
\code{...Failed}; auth-service hoặc xóa hẳn hoặc khôi phục
\code{ACTIVE}. Hai thứ khiến nó đúng thay vì hy vọng: khóa ngữ nghĩa
--- trạng thái \code{DELETING} chặn đăng nhập trong vài giây saga chạy
--- và bước lũy đẳng, vì Kafka sẽ giao ít nhất một lần. Tôi chọn
choreography cho hai hay ba bước và orchestration (Temporal, Camunda,
hoặc một máy trạng thái tự viết) khi đồ thị sự kiện cần một sơ đồ mới
theo dõi được. 2PC bị loại không phải vì nó sai mà vì Kafka và Postgres
không chia sẻ transaction coordinator, và giữ lock xuyên một lời gọi
mạng là cách bạn nhận cuộc gọi lúc 3 giờ sáng.

\subsection*{Q: Một sự kiện bị mất giữa commit cơ sở dữ liệu và lần gửi Kafka. Làm sao để điều đó bất khả?}

Bằng cách không bao giờ gửi từ luồng request. Transactional outbox: sự
kiện là một dòng trong bảng \code{outbox} ghi trong cùng transaction với
dòng nghiệp vụ, nên Postgres bảo đảm cả hai hoặc không gì. Một relay
--- một poller \code{@Scheduled} với \code{FOR UPDATE SKIP LOCKED}, hoặc
Debezium đọc WAL --- phát nó sau đó và đánh dấu. Relay là
ít-nhất-một-lần, nên dòng mang một id sự kiện để consumer khử trùng.
Trong dự án này nó còn gỡ một sự cố tiềm ẩn: hôm nay
\code{kafkaTemplate.send()} bên trong \code{AuthService.register} chặn
trong \code{max.block.ms} khi broker chết, nên một sự cố broker mail
khiến đăng ký thất bại sau sáu mươi giây. Với outbox, đăng ký commit
trong mili giây và e-mail đến trễ thay vì lượt đăng ký bị mất.

\subsection*{Q: Làm sao chuyển một cơ sở dữ liệu dùng chung sang mỗi service một cơ sở dữ liệu một cách an toàn?}

Strangler, không phải big bang. Bước một: cầm máu --- tìm mọi join và
khóa ngoại xuyên schema và định tuyến chúng qua API của service sở hữu
hoặc một bản sao cục bộ, trong khi các bảng vẫn ở một cơ sở dữ liệu. Đó
là phần khó và có thể mất nhiều tháng; tách vật lý sau đó là một cuối
tuần. Bước hai: cấp cho mỗi service credential và schema riêng, để một
join vô tình thất bại ngay lúc biên dịch. Bước ba: ghi kép hoặc CDC các
bảng sang cơ sở dữ liệu mới, kiểm đếm và checksum, lật phần đọc, rồi
phần ghi, giữ bảng cũ chỉ-đọc thêm một bản phát hành, rồi drop. Xuyên
suốt, outbox đi vào \emph{trước} khi tách, vì khoảnh khắc join biến mất
là bạn dựa vào sự kiện để giữ bản sao đúng. ``Hai server Postgres,
không phải hai schema'' của Chương~7 là trạng thái cuối; bạn tới đó
bằng schema-mỗi-service trước.

\subsection*{Q: Bạn đã nhân bản dữ liệu Student vào course-service. Làm sao giữ nó đúng?}

Tôi chấp nhận nó sẽ cũ và tôi chặn trên độ cũ. Việc tạo tới trên
\code{ts.user.activated}; cập nhật cần một \code{ts.user.updated} dự án
chưa phát --- đó là sửa đầu tiên. Cả hai key theo \code{userId}, nên
chung partition và cập nhật không thể vượt mặt việc tạo. Listener upsert
theo khóa tự nhiên, khiến phát lại an toàn. Sự kiện tới từ outbox, nên
không cái nào mất; consumer ghi lại id sự kiện đã xử lý, nên không cái
nào áp dụng hai lần. Rồi tôi đo: một job hằng đêm so một mẫu
\code{student.email} với auth-service và cảnh báo khi trôi, vì mỗi cơ
chế ở trên đều có một lỗi tôi chưa tìm ra. Điều tôi không làm là đặt
e-mail vào JWT --- việc đó đóng băng nó suốt đời token và ghép schema
qua hợp đồng token (Chương~28).

\begin{quiz}
\begin{enumerate}
  \item Đọc \code{AuthService.activate}. Mô tả hai thứ tự sự cố khác
  nhau khiến auth-service và course-service bất đồng, và nói thứ tự nào
  tạo ra một dòng lẽ ra không tồn tại.
  \item Kafka không với được. Giải thích vì sao \code{POST /register}
  hôm nay mất khoảng một phút để thất bại và không tạo user, nêu tên
  thiết lập producer chịu trách nhiệm. Cùng request đó làm gì sau khi
  có outbox?
  \item Trong relay, vì sao \code{FOR UPDATE SKIP LOCKED} là cần thiết,
  và chuyện gì hỏng với hai replica auth-service nếu thay nó bằng một
  \code{SELECT} thuần?
  \item Vì sao \code{break} (không phải \code{continue}) theo sau một
  lần gửi thất bại trong vòng lặp relay? Nêu kịch bản mà phương án kia
  làm hỏng.
  \item E-mail của một học viên đổi trong auth-service. Nêu
  course-service biết gì về việc đó hôm nay, và liệt kê ba thay đổi
  code khép lại khoảng trống.
  \item Thiết kế xóa tài khoản xuyên auth-service, course-service và
  MinIO như một saga choreography: gọi tên các trạng thái, các sự kiện,
  hành động bù, và giải thích trạng thái \code{DELETING} bảo vệ khỏi
  điều gì.
\end{enumerate}
\end{quiz}
